\item \model{Magistral}

\paragraph*{مقدمه و تبیین ماهیت دادگان مدل}
مدل‌های خانواده \model{Magistral} به عنوان نخستین مدل‌های استدلالی شرکت \en{Mistral AI} \cite{mistral2025magistral} معرفی شده‌اند. برخلاف رویه سنتی پیش‌آموزش خام بر روی تریلیون‌ها توکن وب، این پژوهش مستقیماً از چک‌پوینت‌های پیش‌آموزش‌دیده پایه‌ای شامل \en{Mistral Medium 3 Instruct} و \en{Mistral Small 3 Instruct} بهره جسته است. تمرکز محوری در این مدل، مهندسی و گردآوری داده‌های استدلالی تاییدپذیر (\en{Verifiable Rewards}) برای بهینه‌سازی یادگیری تقویتی (\en{RLVR}) و فاز راه‌اندازی نظارت‌شده (\en{SFT Bootstrapping}) است:
\begin{itemize}
    \item \textbf{مدل \en{Magistral Medium}:} آموزش‌دیده بر بستر \en{Mistral Medium 3} صرفاً از طریق یادگیری تقویتی خالص (\en{Pure RL}) بدون بهره‌گیری از هرگونه تقطیر یا ردپای استدلالی اولیه، که منجر به جهش نزدیک به ۵۰ درصدی در بنچمارک \en{AIME-24} گردیده است.
    \item \textbf{مدل \en{Magistral Small}:} مدلی ۲۴ میلیارد پارامتری بر پایه \en{Mistral Small 3} که فرایند آموزش آن با تلفیق داده‌های آغازین سرد (\en{Cold-start Data}) استخراج‌شده از نسخه \en{Medium} و سپس اجرای برخط \en{RL} تکمیل شده است.
\end{itemize}

\paragraph*{پالایش چندمرحله‌ای دادگان ریاضی (\en{Math Data Curation})}
فرایند گردآوری دادگان ریاضی با هدف دست‌یابی به مسائلی با پاسخ‌های تحلیلی، قطعی و ارزیابی‌پذیر توسط مفسرهای نمادین طراحی شده است:
\begin{itemize}
    \item \textbf{پالایش ساختار و فرمت (\en{Format Filtering}):} کار با یک مجموعه بزرگ اما نویزدار شامل حدود $700\text{K}$ مسئله آغاز گردید. به منظور امکان‌پذیر ساختن اعتبارسنجی قطعی، تمامی مسائل اثباتی تشریحی و مسائل چندبخشی که ارزیابی آنها با قواعد الگوریتمی دشوار بود، حذف شدند. همچنین پرسش‌های چندگزینه‌ای به مسائل گزاره‌ای و تشریحی بازنویسی شدند تا احتمال حدس تصادفی خنثی شده و استخراج پاسخ در قالب ساختار جبری \en{SymPy} درون جعبه \lr{\texttt{\textbackslash boxed\{\}}} مقدور گردد. این گام حجم داده‌ها را به $501\text{K}$ کاهش داد.
    \item \textbf{فیلتر دوسطحی دشواری و ناحیه طلایی (\en{Goldilocks Zone Filtering}):} برای جلوگیری از افت گرادیان ناشی از نمونه‌های بسیار آسان یا ناممکن، یک غربالگری دومرحله‌ای اعمال شد:
    \begin{enumerate}
        \item در مرحله نخست، با نمونه‌برداری ۱۶ پاسخ مستقل به ازای هر مسئله توسط مدل \en{Mistral Large 2} \cite{mistral2024large2}، مسائلی با نرخ موفقیت صفر یا صد درصد کنار گذاشته شدند و داده‌های باقی‌مانده صرف آموزش یک چک‌پوینت نمره‌گذار ۲۴ میلیارد پارامتری شدند.
        \item در مرحله دوم، این مدل استدلالی ۲۴ میلیارد پارامتری کل دادگان را مجدداً ارزیابی نمود. مسائلی که اکثریت خروجی‌های مدل در آنها بر پاسخی یکسان اجماع داشتند اما با برچسب واقعیت زمینی (\en{Ground Truth}) داده اولیه مغایر بودند، به دلیل خطا در برچسب مرجع حذف گردیدند.
    \end{enumerate}
\end{itemize}
در نتیجه این فیلترینگ سخت‌گیرانه، داده‌های نهایی ریاضی به \textbf{۳۸ هزار مسئله فوق‌العاده باکیفیت} محدود شد.

\paragraph*{گردآوری و مهندسی دادگان برنامه‌نویسی (\en{Code Data Curation})}
پیکره کدنویسی رقابتی با هدف سنجش دقیق درستی و کارایی الگوریتم‌ها شکل گرفته است:
\begin{itemize}
    \item \textbf{اعتبارسنجی و تصحیح آزمون‌های واحد (\en{Unit Tests}):} مسائل فاقد آزمون یا کدهای مرجع معتبر حذف شدند. راه‌حل‌های مرجع بر روی آزمون‌ها اجرا شده و آزمون‌های ناهماهنگ پالایش شدند. در مواردی که تمامی راه‌حل‌ها در یک تست با شکست مواجه می‌شدند، برچسب تست بر مبنای خروجی متداول راه‌حل‌ها تصحیح شد. در صورت نیاز، تست‌های مصنوعی جدیدی نیز تولید و اضافه گردید.
    \item \textbf{پشتیبانی دوزبانه و استانداردسازی کامپایل:} مسائل به گونه‌ای بازطراحی یا تکثیر شدند که پیاده‌سازی به هر دو زبان \en{Python} و \en{C++20} پشتیبانی شود. برای تسریع زمان ارزیابی، هدر پرکاربرد \lr{\texttt{bits/stdc++.h}} پیش‌کامپایل گردید.
    \item \textbf{محیط ارزیابی ایمن (\en{Sandbox Execution}):} برای هر ارزیابی گروهی، ۲۰ آزمون تصادفی انتخاب و به تمامی اعضای گروه اختصاص یافت. سقف زمان اجرای هر تست ۴ ثانیه و محدودیت حافظه $300\text{ MB}$ تعیین گردید. این خط لوله در نهایت به گردآوری \textbf{۳۵ هزار مسئله کدنویسی} انجامید.
\end{itemize}

\paragraph*{مهندسی دادگان چندزبانه و پایداری زبان تفکر (\en{Multilingual Strategy})}
به منظور مهار پدیده سوئیچ ناخواسته زبان در طول زنجیره تفکر (\en{Language Mixing}):
\begin{itemize}
    \item \textbf{ترجمه متوازن دادگان:} دقیقاً $10\%$ از مسائل انگلیسی به زبان‌های فرانسوی، اسپانیایی، ایتالیایی، آلمانی، چینی و روسی برگردانده شد.
    \item \textbf{پاداش پایداری زبان (\en{Language Consistency}):} پس از پاک‌سازی علائم لاتک و کدهای دستوری، سه‌گانه پرسش، افکار داخل تگ \lr{\texttt{<think>}} و پاسخ نهایی توسط طبقه‌بند \en{fastText} \cite{joulin2016fasttext} بررسی شده و در صورت یکسانی زبان در هر سه بخش، پاداش افزوده‌ای معادل $0.1$ اعمال گردید. این تمهید امکان تولید زنجیره تفکر به زبان بومی کاربر را بدون افت کیفیت استدلال تثبیت کرد.
\end{itemize}

\paragraph*{دادگان راه‌اندازی اولیه و تقطیر برای مدل کوچک (\en{SFT Cold-Start Bootstrapping})}
در آموزش مدل \model{Magistral Small}، برای آماده‌سازی اولیه سیاست مدل:
\begin{itemize}
    \item \textbf{استخراج ردپاهای باکیفیت آموزگار:} نمونه پاسخ‌های صحیح از مراحل پیشرفته آموزش \en{RL} مدل \en{Magistral Medium} استخراج و ردهای اولیه با تفکر کوتاه فیلتر شدند.
    \item \textbf{غنی‌سازی با داده‌های متن‌باز:} پرومپت‌های متنوع از مجموعه‌های \en{OpenThoughts} \cite{guha2025openthoughts} و \en{OpenR1} \cite{openr1_2025} گزینش شده و با پاسخ‌های تفکری مدل \en{Medium} غنی‌سازی شدند.
    \item \textbf{تزریق داده‌های عمومی دستورپذیری:} برای پیشگیری از فراموشی فاجعه‌بار مهارت‌های عمومی، $10\%$ از حجم داده‌های \en{SFT} به دستورات متداول غیر‌استدلالی اختصاص یافت و مدل به مدت ۴ دوره بر روی این ترکیب آموزش دید.
\end{itemize}

\paragraph*{صورت‌بندی بهینه‌سازی و سیگنال‌های پاداش حاکم بر دادگان}
فرایند بهینه‌سازی سیاست مدل $\pi_\theta$ بر روی دادگان پردازش‌شده توسط الگوریتم \en{GRPO} \cite{shao2024deepseekmath} و با اعمال اصلاحات تثبیت‌کننده به فرم زیر اجرا می‌شود:
\begin{equation}
    \mathcal{J}_{\text{GRPO}}(\theta) = \mathbb{E}_{q \sim \mathcal{P}(Q), \{o_i\}_{i=1}^G \sim \pi_{\theta_{\text{old}}}} \left[ \frac{1}{\sum_{i=1}^G |o_i|} \sum_{i=1}^G \sum_{t=1}^{|o_i|} \min \left( r_{i,t}(\theta) \hat{A}_{i,t}^{\text{norm}}, \text{clip}(r_{i,t}(\theta), 1 - \varepsilon_{\text{low}}, 1 + \varepsilon_{\text{high}}) \hat{A}_{i,t}^{\text{norm}} \right) \right]
\end{equation}
که در آن $r_{i,t}(\theta) = \frac{\pi_\theta(o_{i,t} \mid q, o_{i,<t})}{\pi_{\theta_{\text{old}}}(o_{i,t} \mid q, o_{i,<t})}$، نسبت مزیت استانداردشده در سطح مینی‌بچ به صورت $\hat{A}_{i,t}^{\text{norm}} = \frac{\hat{A}_i - \hat{A}_{\text{mean}}}{\hat{A}_{\text{std}}}$ تعریف شده و شرط تفکیک گروه‌های همگن $\exists m < n, r_m \neq r_n$ اعمال می‌شود.

همچنین جهت پیشگیری از رشد مهارگسیخته طول پاسخ و سرریز حافظه، جریمه طول نرم بر داده‌ها حاکم گردید:
\begin{equation}
    R_{\text{length}}(y) = \begin{cases} 
        0, & |y| \le l_{\text{max}} - l_{\text{cache}} \\ 
        -0.1 \cdot \frac{|y| - l_{\text{max}} + l_{\text{cache}}}{l_{\text{cache}}}, & l_{\text{max}} - l_{\text{cache}} < |y| \le l_{\text{max}} \\ 
        -0.1, & l_{\text{max}} < |y| 
    \end{cases}
\end{equation}

\paragraph*{تحلیل‌های آماری، پویایی وزن‌ها و ارزیابی جامع دادگان}
داده‌های آموزشی و اثرات آنها بر روی شاخص‌های کمی در جداول زیر خلاصه شده است. جدول \ref{tab:magistral_math_filtering} فرآیند فشرده‌سازی و غربالگری دادگان ریاضی را نشان می‌دهد که در آن تنها ۵.۴٪ از داده‌های اولیه موفق به عبور از فیلترهای صحت و دشواری شدند. جدول \ref{tab:magistral_benchmarks} نیز نشان‌دهنده ارتقای چشمگیر عملکرد استدلالی مدل در حوزه‌های گوناگون تحت آموزش بر روی این دادگان گزینش‌شده است.

\begin{table}[htbp]
\centering
\caption{مراحل متوالی پالایش کمی دادگان ریاضی در مدل‌های \model{Magistral}}
\label{tab:magistral_math_filtering}
\begin{tabular}{lccc}
\toprule
\textbf{شاخص} & \textbf{داده اولیه خام (\en{Initial})} & \textbf{پالایش فرمت و ساختار} & \textbf{پالایش دشواری ناحیه طلایی} \\
\midrule
تعداد نمونه‌ها & $699\text{K}$ & $501\text{K}$ & $38\text{K}$ \\
نرخ بقا نسبت به مرحله قبل & - & $71.67\%$ & $7.58\%$ \\
نرخ تجمعی بقا از کل & $100\%$ & $71.67\%$ & \textbf{5.44\%} \\
\bottomrule
\end{tabular}
\end{table}

\begin{table}[htbp]
\centering
\caption{ارزیابی عملکرد مدل‌های \model{Magistral} در مقایسه با مبناهای استدلالی تراز اول جهان}
\label{tab:magistral_benchmarks}
\resizebox{\textwidth}{!}{%
\begin{tabular}{lcccccc}
\toprule
\textbf{بنچمارک (متریک)} & \textbf{\en{Mistral Med. 3}} & \textbf{\en{Magistral Med.}} & \textbf{\en{DeepSeek-V3}} & \textbf{\en{DeepSeek-R1-Zero}} & \textbf{\en{DeepSeek-R1}} \\
\midrule
استفاده از \en{SFT} قبل از \en{RL} & - & \textbf{خیر} & - & خیر & بله \\
\midrule
\en{AIME '24 (Pass@1 / Maj@64)} & 26.8 / 43.4 & \textbf{73.6 / 90.0} & 39.2 & 71.0 / - & 79.8 / - \\
\en{AIME '25 (Pass@1 / Maj@64)} & 21.2 / 30.0 & \textbf{64.9 / 83.3} & 28.8 & - & 70.0 / - \\
\en{MATH-500} & 91.0 & \textbf{94.3} & 90.2 & 95.9 & 97.3 \\
\en{GPQA Diamond} & 59.6 & \textbf{70.8} & 59.1 & 73.3 & 71.5 \\
\en{LiveCodeBench v5} & 29.1 & \textbf{59.4} & 36.2 & 50.0 & 65.9 \\
\en{Aider Polyglot} & 28.9 & \textbf{47.1} & 49.6 & - & 53.3 \\
\en{Humanity's Last Exam (Text)} & 4.4 & \textbf{9.0} & - & - & 8.6 \\
\bottomrule
\end{tabular}%
}
\end{table}

\paragraph*{یافته‌های انتقال بدون داده و راهبردهای ناموفق (\en{Negative Results})}
تحلیل پویایی آموزش مدل حاوی دو بینش اساسی در حوزه مهندسی دادگان است:
\begin{itemize}
    \item \textbf{انتقال چندوجهی بدون دادگان تصویر (\en{Multimodal Free Lunch}):} علیرغم اینکه آموزش \en{RL} کاملاً بر بستر داده‌های متنی محض صورت پذیرفت، مدل قابلیت‌های انکودر بصری پیش‌آموزش‌دیده خود را حفظ کرده و حتی در بنچمارک‌های چندوجهی ارتقا یافت؛ به گونه‌ای که امتیاز \en{MMMU} از $65.0\%$ به $70.0\%$ و در \en{MMMU-Pro (Vision)} با رشدی چشمگیر از $39.7\%$ به $52.1\%$ رسید.
    \item \textbf{شکست پاداش جزئی در دادگان کد:} فرضیه اعطای پاداش نسبی بر پایه درصد قبولی آزمون‌های واحد به جای پاداش دودویی ($0$ یا $1$)، به افت ۲ درصدی دقت در \en{LiveCodeBench} انجامید، زیرا پاداش جزئی منجر به تقویت الگوهای غیربهینه و ایجاد سیگنال گمراه‌کننده در کدهای دارای باگ می‌شد.
\end{itemize}